\section{The AI--human Analogy Is Structurally Asymmetric}
\label{sec:asymmetries}

The matched examples in this paper should not be read as claiming that AI optimization
is simply accelerated human problem solving.  The comparison is useful because the
same feedback structures recur in both cases, but several important mechanisms are
asymmetric.

\paragraph{AI has much greater parallel search.}
A human expert normally explores a small number of candidate solutions at once.
Digital systems can generate, execute, and reject enormous candidate sets in parallel.
AlphaProof trains on millions of formally checked problems
\parencite{hubert2025alphaproof}; AlphaDev explored millions of candidate programs
\parencite{mankowitz2023alphadev}.  Consequently, wall-clock ratios such as
\(\tauT/\tauP\) can substantially understate the amount of selection occurring
between target checks.  This is one reason we define \(\Nproxy\) primarily as the number of
consequential proxy-guided updates, rather than as a time ratio.

\paragraph{AI state can be copied and branched cheaply.}
Human learning is transmitted imperfectly through teaching, writing, imitation, and
institutions.  A trained digital state can instead be copied, checkpointed, branched,
evaluated, and recombined with comparatively little loss.  A useful update can
therefore propagate across many simultaneous search processes without requiring each
process to rediscover it.  There is no close human analogue at the level of individual
cognition.

\paragraph{AI can be optimized internally.}
Human feedback primarily selects actions, strategies, institutions, or people.
Machine learning also changes the internal parameters that generate subsequent
behavior.  Gradient-based optimization can therefore perform many small internal
updates between externally meaningful decisions.  Whether each gradient step should
count toward \(\Nproxy\) depends on the analysis.  For safety questions, the relevant unit
is normally the smallest update capable of producing a consequential change in the
system's policy or deployment behavior, not every floating-point parameter update.

\paragraph{Human progress has stronger examples of verifier invention.}
The asymmetry also runs in the other direction.  Human scientific progress repeatedly
changed not only candidate solutions but the \emph{measurement regime}: new
instruments, standardized measurements, controlled experiments, statistical methods,
and institutions made previously weakly observable claims testable.  Current AI
systems can propose experiments and help design measurement procedures, and
self-driving laboratories automate existing instances
\parencite{volk2023alphaflow,szymanski2023alab}.  But the historical record contains much
stronger examples of humans changing what could be observed at all.

\paragraph{Humans can sometimes learn from \(N=1\) events.}
Wars, accidents, financial crises, constitutional failures, and scientific anomalies
may occur too rarely for ordinary repeated optimization.  Humans nevertheless build
causal narratives, compare mechanisms across cases, construct counterfactuals, and
change institutions after unique events.  These updates can be wrong, but they are not
well represented as repeated trials from a stationary distribution.  Present AI
training, by contrast, derives much of its strength from large repeated datasets or
synthetically generated task families.

\paragraph{Human institutions can deliberately separate evaluators.}
Science, auditing, courts, safety engineering, and other institutions often create
partially independent channels of criticism.  Their purpose is not merely to increase
feedback frequency but to prevent one selection criterion from controlling all
evidence about itself.  AI evaluation often compresses many judgments into a single
benchmark or reward model.  Multiple AI critics do not automatically recreate
institutional independence if they share training data, architecture, prompts, or
failure modes.

\paragraph{The comparison is therefore conditional.}
The paper makes a comparison at the level of feedback architecture:
\[
    \text{candidate}
    \rightarrow
    \text{selection signal}
    \rightarrow
    \text{update}
    \rightarrow
    \text{later target evidence}.
\]
It does not claim equivalence between human cognition, scientific institutions, and
machine learning.  The important empirical question is narrower: holding the target
and evaluator structure approximately fixed, what changes when the number of
proxy-guided selections between independent target checks increases?

This scope restriction matters for interpreting the historical examples.  The human
record motivates the variables and provides useful analogies, but it does not by
itself establish a causal law for AI systems.  Conversely, reward-model
overoptimization provides direct evidence for the accumulation mechanism in AI but
does not establish that institutional or scientific systems obey the same quantitative
relationship.
